% ============================================================================
%  SECTION 4 — OUTILS THEORIQUES : SOUS-DIFFERENTIEL & OPERATEUR PROXIMAL
%  >>> A REMPLIR PAR LE MEMBRE 2 <<<
%
%  Objectif de la section : fournir la fondation theorique dont dependent les
%  preuves des sections 5 (ISTA) et 6 (FISTA). C'est le PIVOT du rapport.
%
%  Contenu attendu (voir REPARTITION_DETAILLEE.md) :
%    - Sous-differentiel d'une fonction convexe non lisse ; cas de |t| en 0
%      (l'ensemble des pentes [-1, 1]).
%    - Condition d'optimalite : 0 \in \partial F(x*).
%    - Definition de l'operateur proximal :
%        prox_{t g}(z) = argmin_u  g(u) + (1/2t) ||u - z||^2
%      + interpretation geometrique.
%    - DERIVATION COMPLETE du seuillage doux (prox de lambda||.||_1) :
%      separation coordonnee par coordonnee -> sign(z)*max(|z|-t*lambda, 0).
%    - Propriete de non-expansivite du prox (utile pour les preuves de conv.).
%
%  Notations a respecter (definies dans main.tex) :
%    \prox, \softthr, \norm{.}, \argmin, \sign, \xopt, \Fopt, \R, \T
%  Reutiliser : F = f + g, f(x)=1/2||Ax-y||^2, g(x)=lambda||x||_1, L=||A^T A||.
% ============================================================================
\section{Outils theoriques : sous-differentiel et operateur proximal}
\label{sec:theorie}

% --- [MEMBRE 2] Ecrire la section ici. Squelette suggere ci-dessous. --------

\subsection{Le sous-differentiel}
% TODO (Membre 2)

\subsection{Condition d'optimalite}
% TODO (Membre 2)

\subsection{L'operateur proximal}
% TODO (Membre 2)

\subsection{Derivation du seuillage doux}
% TODO (Membre 2) : c'est le coeur. La formule a etablir est
% \prox_{t\lambda\norm{\cdot}_1}(z)_i = \sign(z_i)\max(|z_i| - t\lambda, 0).

\subsection{Non-expansivite}
% TODO (Membre 2)
